\clearpage
\thispagestyle{empty}
\pdfbookmark[0]{Beobachtungen der magnetischen Inklination in Göttingen. Ausgabe und Quellen}{inclination1841-edition}
\begin{center}\Large\textbf{Beobachtungen der magnetischen Inklination in Göttingen}\end{center}

Carl Friedrich Gauss’ Abhandlung erschien in den \textit{Resultaten aus den Beobachtungen des magnetischen Vereins} für 1841 als Artikel II. Die KI-Übertragung aller 34 Artikel und sechs historischen Anmerkungen in modernes Deutsch mit KI-erstelltem Satz folgt dem deutschen Text der \textit{Werke}, Bd.~V (Göttingen, 1867), S.~444--492. Die Beobachtungen, Tabellen und die Schreibweise der Zahlen folgen dieser Ausgabe.

\medskip\noindent
\href{https://gdz.sub.uni-goettingen.de/id/PPN236006339}{Gesammelte deutsche Ausgabe (1867)}\par
\noindent\href{https://doi.org/10.1007/978-3-642-49319-5_22}{Deutscher Nachdruck (1877)}\par
\noindent\href{https://collections.whipplemuseum.cam.ac.uk/objects/12342/}{Verwandtes Robinson-Inklinometer, Whipple Museum, Inventarnummer 3781}

\medskip\noindent
Zu den von Gauss zitierten Quellen gehören Johann Caspar Horners \href{https://api.digitale-sammlungen.de/iiif/image/v2/bsb10131088_00105/full/1500,/0/default.jpg}{Artikel über das Inklinatorium}, der Bericht von Humphrey Lloyd, Edward Sabine und James Clarke Ross über \href{https://archive.org/download/reportofbritisha05unse/page/n207.jpg}{magnetische Beobachtungen in Irland} und Adolph Ermans \href{https://resolver.sub.uni-goettingen.de/purl?PPN329604732}{physikalische Beobachtungen}. Der abschließende Vergleich verweist auf \href{https://books.google.com/books?id=O_kOAAAAQAAJ&pg=PA152}{Alexandre de Humboldts Göttinger Angaben}, \href{https://doi.org/10.1017/S008045680002980X}{James D. Forbes’ Messungen} und \href{https://gdz.sub.uni-goettingen.de/id/PPN35283028X_0003_2NS}{Johann Tobias Mayers Untersuchung der Inklinationsnadel}.
\clearpage
